\songtitle{Villegas de Scholer}

\tag*{2022}\tag{Scholer}\tag{victoria-eugenia}
\begin{notes}
Lyrics by Claude and Carlos Thompson, based on character Victoria Eugenia Villegas de Scholer, %
from the character dossier document for the planed series \emph{Scholer} %
written and conceptualized by Carlos Thompson.
\end{notes}
\begin{style}
latin romantic ballad, cinematic chamber pop, 1980s ballad, female mezzo vocals, gang choir chorus, nylon-string tremolo, piano arpeggios, string quartet lines, oboe countermelody, muted trumpet, drum machine backbeat, brushed snare, gated reverb, tape saturation, plate reverb, stereo delay throws, dramatic crescendo, bittersweet scandal
\end{style}
\begin{lyrics}
\songpart{Verse 1}
De Cisneros viene, de estirpe antioqueña,\\
los Villegas Ángel, ganado y hacienda,\\
apellido que abre las puertas de piedra,\\
alta sociedad paisa, pedigrí y leyenda.

\songpart{Verse 2}
Se conocieron los dos en una feria ganadera,\\
dos apellidos que unieron su bandera,\\
se casaron y llegaron juntos a Bogotá,\\
donde empezaron su casa y su manera.

\songpart{Chorus}
Victoria Eugenia Villegas de Scholer,\\
la dama de sociedad, la dama sin doblez,\\
matrimonio que el mundo mira con respeto,\\
secretos que la ciudad no debe conocer.

\songpart{Verse 3}
Agustín ya se forma en ingeniería,\\
María Luisa llegó de Nueva York un día,\\
Andrea Sofía se gradúa con alegría,\\
y Nicolás, el menor, inicia su travesía.

\songpart{Bridge}
Ella también camina por sendas propias,\\
Roberto entre sus amantes, discreto en las sombras,\\
dos infieles que se cruzan sin reproche,\\
cada uno con su vida, cada cual con su otra.

\songpart{Final Chorus}
Victoria Eugenia Villegas de Scholer,\\
la dama de sociedad, la dama sin doblez,\\
matrimonio que el mundo mira con respeto,\\
secretos que la ciudad no debe conocer.\\
Que se sepa esto, y el mundo entero se entera,\\
se hunde el nombre, se hunde la casa,\\
no hay quien la sostenga después.\\
no se sostiene después.
\end{lyrics}

